\chapter{消失猜想}\label{ch:9}
\begingroup\setlength{\belowdisplayshortskip}{2pt}
\titlespacing*{\section}{0pt}{1.5ex plus .5ex}{1ex plus .2ex}
\srcpage{48}
本章解释\chref{6}洁净延拓论证和\secref{8.6}尖点性论证所用的平均消失定理\cite{Gai04,Gai16b}。记号与此前相同：$X$ 是亏格为 $g$ 的光滑射影曲线，$D(-)$ 表示所选层理论。以下 de Rham 讨论取 $k=\mathbb C$、$D(-)=\DMod(-)$。
\section{平均函子}\label{sec:9.1}
设 $E$ 为秩 $r$ 的不可约局部系统，$m<r$。对 $d=d_2-d_1\ge0$，考虑秩 $m$ 丛的上修改：
\origdiagram{p048-diagram.pdf}{0.3962}
其中 $\operatorname{Mod}_m^{d_1,d_2}$ 的一点为 $M\subset M'$，$\length(M'/M)=d$，且 $\pi(M\subset M')=M'/M$。令 $\cL_E^d\in D(\Tor^d)$ 为 Laumon 层：在支撑映射 $\Tor^d\to X^{(d)}$ 下，无重数轨迹 $X^{(d),\circ}\subset X^{(d)}$ 的逆像上，它是 $E^{\boxtimes d}$ 的对称下降 $E^{(d)}$；整体上，它是后者到 $\Tor^d$ 的 Goresky–MacPherson 延拓。忽略通常的归一化移位与 Tate 扭转，
\[\Av_E^d(\cF):=\overrightarrow h_!(\overleftarrow h^*\cF\otimes\pi^*\cL_E^d)\in D(\Bun_m^{d_2}).\]
等价地，对各连通分支求和后，这是 $D(\Bun_m)$ 的自函子。称其为“由 $\cL_E^d$ 作上修改”；文献\cite{Gai16b}从谱一侧回顾同一个平均函子。
\begin{theorem}{消失定理}{9.1.1}
假设 $E$ 不可约、$\rank(E)=r>m$，并为简便假设 $g>1$。则只要 $d>rm(2g-2)$，就有
\[\Av_E^d=0.\]
采用移位与扭曲的标准归一化，同样的表述适用于通常的层论情形。
\end{theorem}
\thref{9.1.1}就是此前以 $r=n$、$m=k<n$ 使用的定理。这里的重点不是构造 $\cL_E^d$，而是从几何 Langlands 的谱一侧理解这个函子为何应当消失。
\section{谱作用的启发式解释}\label{sec:9.2}\index{谱作用|hyperpage}\index{温和对象|hyperpage}
在几何 Langlands 的 de Rham 形式中，纲要上预期有文献\cite{GR24a,AG15}的谱作用及幂零支撑表述：
\[\DMod(\Bun_G)^{\mathrm{temp}}\simeq\QCoh(\LocSys_{\checkG}),\qquad\DMod(\Bun_G)\simeq\IndCoh_{\mathrm{Nilp}}(\LocSys_{\checkG}).\]
将右边限制到紧对象，这常常启发性地表现为谱叠上的凝聚层。因此，$\QCoh(\LocSys_{\checkG})$ 应当作用于 $\DMod(\Bun_G)$。
这一作用并非任意的。对 $E\in\LocSys_{\checkG}$，纤维函子
\[(-)_E:\QCoh(\LocSys_{\checkG})\longrightarrow\operatorname{Vect}\]
应通过基变换切出 $E$-Hecke 范畴：
\[\DMod(\Bun_G)\otimes_{\QCoh(\LocSys_{\checkG})}\operatorname{Vect}\simeq\DMod(\Bun_G)_{E\text{-Hecke}}.\]
\srcpage{49}
这里，右边由 Hecke 变换可通过与局部系统 $E$ 作张量积得到的对象组成。

在固定点 $x\in X$ 处，几何 Satake 给出函子
\[\Rep(\checkG)\otimes D(\Bun_G)\longrightarrow D(\Bun_G),\qquad(V,\cK)\longmapsto H_{V,x}(\cK),\]
其中 $\Rep(\checkG)$ 是球面范畴 $D(G(\cO_x)\backslash G(K_x)/G(\cO_x))$ 内的 Satake 心。若让点移动，则得到
\[\Rep(\checkG)\otimes D(\Bun_G)\longrightarrow D(\Bun_G\times X).\]
$E$-Hecke 特征层是配备同构
\[H_V(\cK)\simeq\cK\boxtimes V_E\]
的对象 $\cK$，其中同构对所有 $V\in\Rep(\checkG)$ 给出，并配备在 $X^I$ 上的高阶相容性。

朴素函子 $\Rep(\checkG)\otimes D(X)\to\End(D(\Bun_G))$ 虽然是函子，却不是所需意义下的幺半函子：当两个点碰撞时，卷积由几何 Satake 的幺半结构控制，而不是由 $X^2$ 上的外张量积控制。例如，必须有
\begin{itemize}
\item 当 $x\ne y$ 时，$H_{V,x}\circ H_{W,y}=H_{V,W;x,y}$；
\item $H_{V,W;x,x}=H_{V\otimes W,x}$。
\end{itemize}
这恰好是因子化条件，所以正确的参数空间必须是 Ran 空间。
\section{Ran 作用与局部到整体}\label{sec:9.3}\index{Ran 空间|hyperpage}\index{因子化范畴|hyperpage}
令
\[\operatorname{Ran}(X)=\{I\subset X\mid0<|I|<\infty\}.\]
它理解为各幂 $X^I$ 沿满射所形成的余极限。因子化范畴 $\Rep(\checkG)_{\mathrm{Ran}}$ 在有限集合 $I\subset X$ 上的纤维为 $\bigotimes_{x\in I}\Rep(\checkG)$，点碰撞时进行融合。Hecke 函子组合成幺半作用

\begin{transnote}{含糊点拨}
取「沿满射」的余极限，是因为 $X^I$ 只记住带标号的点，而满射 $I\twoheadrightarrow J$ 允许把若干点合并；把这些粘起来才得到「任意有限子集都能塞进去、点碰撞时发生融合」的 Ran 空间。纤维写成 $\bigotimes_{x\in I}$ 而不是 $\prod$：$\Rep(\check G)$ 是张量范畴，纤维里的对象是各点处表示的张量积，融合把两点的张量积在碰撞极限下乘成一个。缺了这层，定理 9.3.2 的幺半作用无从谈起。[来源：nLab "Ran space"，https://ncatlab.org/nlab/show/Ran+space（Last revised 2020-07-27；2026-11-17 访问）；并见 nLab "factorization algebra"（Last revised 2026-08-16）。]
\end{transnote}
\[\Rep(\checkG)_{\mathrm{Ran}}\longrightarrow\End(\DMod(\Bun_G)).\]
还有幺半“局部化”函子
\[\operatorname{Loc}_{\checkG}:\Rep(\checkG)_{\mathrm{Ran}}\longrightarrow\QCoh(\LocSys_{\checkG}),\]
它将集合 $(V_x)_{x\in I}$ 映为 $\LocSys_{\checkG}$ 上的取值向量丛 $\bigotimes_{x\in I}V_{E,x}$。

以下结果见文献\cite{Gai10}。
\begin{theorem}{Gaitsgory–Lurie，局部到整体}{9.3.1}
局部化函子 $\operatorname{Loc}_{\checkG}$ 将 $\QCoh(\LocSys_{\checkG})$ 实现为 $\Rep(\checkG)_{\mathrm{Ran}}$ 的 Verdier 商。等价地，$\operatorname{Loc}_{\checkG}$ 的右伴随全忠实。
\end{theorem}

证明在形式上来自描述
\[\LocSys_{\checkG}=\{B\checkG\times X\to X\text{ 的水平截面}\};\]
同一论证适用于任何配备相对于 $X$ 的联络的叠 $Z\to X$。因此，知道 Ran Hecke 作用后，谱作用不再需要额外选择。
\srcpage{50}
\begin{theorem}{广义消失定理\cite{Gai10,Gai04}}{9.3.2}
幺半作用
\[\Rep(\checkG)_{\mathrm{Ran}}\longrightarrow\End(\DMod(\Bun_G))\]
唯一地经由 $\QCoh(\LocSys_{\checkG})$ 分解：
\origdiagram{p050-triangle.pdf}{0.4668}
这是广义消失猜想的范畴形式。其证明使用 Kac–Moody 表示的 Beilinson–Drinfeld 局部化。
\end{theorem}
\section{恢复函子 \texorpdfstring{$\Av_E^d$}{Av（E，d）}}\label{sec:9.4}
取 $G=\GL_m$，从而 $\checkG=\GL_m$，令 $\mathrm{St}_m$ 为标准表示。$\Rep(\GL_m)_{\mathrm{Ran}}$ 中与平均核对应的对象为
\[\cA_{E,d}:=(\Sym^d(\mathrm{St}_m\otimes E))_{X^{(d)}}.\]
它支撑于自然映射 $X^{(d)}\to\operatorname{Ran}(X)$。采用与几何 Satake 对应相同的归一化，Ran Hecke 作用将 $\cA_{E,d}$ 映为 $\Av_E^d$。

原因在局部 Satake 纤维上可见。在点 $x\in X$ 上，Cauchy 分解给出
\[\Sym^d(\mathrm{St}_m\otimes E_x)\simeq\bigoplus_{\substack{\mu\vdash d\\\ell(\mu)\le\min(m,r)}}S_\mu(\mathrm{St}_m)\otimes S_\mu(E_x).\]

\begin{transnote}{含糊点拨}
$S_\mu$ 是 Schur 函子，$\mu\vdash d$ 对应 $S_d$ 的不可约表示；Cauchy 分解把 $\Sym^d(V\otimes W)$ 按对称群的不可约分量拆开，所以右端是分拆求和而不是单一张量。截断 $\ell(\mu)\le m$ 来自 $\mathrm{St}_m$ 只有 $m$ 维：$\wedge^i\mathrm{St}_m=0$ 当 $i>m$，长度更长的分拆没有贡献；$\ell(\mu)\le r$ 同理来自 $\rank E_x=r$。这正是把 $\Sym^d$ 化成有限和、可逐项计算的依据。[依据：原书 p.50 本式与 §9.4 关于 Cauchy 分解的说明。]
\end{transnote}
在几何 Satake 下，$S_\mu(\mathrm{St}_m)$ 对应于由 $\mu$ 标记的 Schubert 分层的 IC 层，而 $S_\mu(E_x)=E_x^\mu$。等价地，右边是不变空间 $(\mathrm{St}_m^{\otimes d}\otimes E_x^{\otimes d})^{S_d}$，并不是对每个直和项再取一次 $S_\mu$ 不变量。这是修改核中出现的 Laumon 层 $\cL_E^d$ 的 Satake 含义。

由\thref{9.3.2}，$\Av_E^d$ 从 $\operatorname{Loc}_{\GL_m}(\cA_{E,d})$ 得到。若 $L$ 为 $X$ 上的局部系统，记 $L^{(d)}$ 为它到 $X^{(d)}$ 的对称下降。忽略对偶约定的差异，该拟凝聚层在秩 $m$ 局部系统 $E'\in\LocSys_{\GL_m}$ 处的纤维为 $\RGamma(X^{(d)},(E'\otimes E)^{(d)})$。因此，只需证明以下初等消失。
\begin{lemma}{}{9.4.1}
设 $E$ 不可约且秩 $r>m$，$E'$ 是任意秩 $m$ 局部系统。若 $d>rm(2g-2)$，则
\[\RGamma(X^{(d)},(E'\otimes E)^{(d)})=0.\]
\end{lemma}
\begin{proof}
置 $L=E'\otimes E$。由于 $E$ 不可约且 $\rank(E)>\rank(E')$，任意映射 $(E')^\vee\to E$ 都为零；故 $H^0(X,L)=0$。对 $E^\vee$ 应用同一论证得到 $H^0(X,L^\vee)=0$，Verdier 对偶给出 $H^2(X,L)=0$。因此 $\RGamma(X,L)$ 集中在次数 $1$，且
\[\dim H^1(X,L)=(2g-2)\rank(L)=rm(2g-2).\]
Künneth 公式给出
\[\RGamma(X^{(d)},L^{(d)})\simeq\Sym^d\RGamma(X,L)\simeq\bigwedge^dH^1(X,L)[-d].\]
当 $d>\dim H^1(X,L)$ 时，最后的外幂为零。
\end{proof}

于是，谱作用将 $\cA_{E,d}$ 映为零，这就在 de Rham 情形给出了 $\Av_E^d$ 的消失定理。
\section{其他层理论与一般约化群}\label{sec:9.5}\index{受限局部系统叠|hyperpage}
\srcpage{51}
对于 de Rham $D$ 模以外的层理论，前述论证中的谱叠要替换为受限局部系统叠\cite{AGK20a}，非正式地记为 $\LocSys_{\checkG}^{\mathrm{restr}}$。可以用点函子描述它：
\[\LocSys_{\checkG}^{\mathrm{restr}}(S)=\{\Rep(\checkG)\xrightarrow{\otimes}\QCoh(S)\otimes\operatorname{QLisse}(X)\},\]
而在 de Rham 情形，将 $\operatorname{QLisse}(X)$ 替换为 $\DMod(X)$ 就有平行的表达式。在 Betti／可构造情形，已知的因子化陈述落在具有幂零奇异支撑的范畴中，纲要上为
\[\Rep(\checkG)_{\mathrm{Ran}}^{\operatorname{QLisse}}\longrightarrow\QCoh(\LocSys_{\checkG}^{\mathrm{restr}})\longrightarrow\End(D_{\mathrm{Nilp}}(\Bun_G)),\]
因此，它本身并不蕴涵整个 $D(\Bun_G)$ 上完整的消失定理。

对一般约化群，预期陈述相同。设 $V\in\Rep(\checkG)$，$E$ 为不可约向量局部系统，$\rank(E)>\dim V$。当
\[d>(2g-2)\dim(V\otimes E)\]
时，Ran 表示范畴中的对象 $\Sym^d(V\otimes E)$ 应在 $D(\Bun_G)$ 上作用为零。在 de Rham 理论中，由与上文相同的上同调计算，这从广义消失定理推出。在 $\ell$-进情形，这一表述中的相应陈述仍作为猜想性输入处理。
\section{Gaitsgory 最初的消失证明}\label{sec:9.6}
本章还概述 Gaitsgory 关于消失定理的最初证明\cite{Gai04}。论证是几何的，但对应于解析陈述：相关 Rankin–Selberg $L$ 函数是次数有界的多项式。

首先，定义下修改和函子 $\Av_{E^*}^{-d}$。忽略标准移位，$\Av_{E^*}^{-d}$ 同时是 $\Av_E^d$ 的左伴随和右伴随，Verdier 对偶给出

\begin{transnote}{含糊点拨}
下修改函子用 $E^*$ 与负次数构造，与 $\Av_E^d$ 互为伴随：忽略标准移位后它同时是左、右伴随。于是 Hom 里的 $\Av_E^d$ 可以左右挪动，$\Av_E^d=0$ 等价于 $\Hom(\Av_E^d(-),-)=0$；Verdier 对偶把紧支撑版本与 $*$-版本互换，就得到后半式 $\mathbb D\circ\Av_{E^*}^{-d}\simeq\Av_E^d\circ\mathbb D$。书中随后正是用这一条把 $F=0$ 从 Hom 计算里推出，而不是硬算核。[依据：原书 p.51 本段与随后的 Hom 计算；并见 nLab "Verdier duality" 的伴随刻画（该条目无 Last revised，2026-11-17 访问）。]
\end{transnote}
\[\mathbb D\circ\Av_{E^*}^{-d}\simeq\Av_E^d\circ\mathbb D.\]
要证明 $\Av_E^d=0$，只需证明
\[\Av_E^d\circ\Av_{E^*}^{-d}=0.\]
事实上，若 $G:=\Av_{E^*}^{-d}$ 是 $F:=\Av_E^d$ 的左伴随且 $FG=0$，则 $\Hom(G\cK,G\cK)\simeq\Hom(\cK,FG\cK)=0$，所以 $G=0$；由于 $G$ 也是 $F$ 的右伴随，便迫使 $F=0$。对消失范围内的 $d$，一个归纳引理说明 $\Av_E^d(\cF)$ 和 $\Av_{E^*}^{-d}(\cF)$ 是尖点的。因此，只要 $\Av_E^d$ 消去尖点层，所需复合就为零。

其次，证明 Euler 示性数归约所需的正合性。更精确地，首先证明 $\Av_E^d$ 在转到下文描述的辅助 Verdier 商后正合；由于其像是尖点的，且该商在尖点对象上保守，这给出消失论证所需的正合性。然后，标准的 Euler 示性数论证将消失归约为
\[\chi(\Av_E^d(\cF))=0.\]
对 Euler 示性数而言，可以将局部系统 $E$ 替换为另一个同秩局部系统；于是归约到平凡局部系统，尽管正合性输入本身使用了 $E$ 的不可约性。
\srcpage{52}

对平凡秩 $r$ 系统 $\mathbf1^{\oplus r}$，Braverman 的计算给出
\[\Av_{\mathbf1^{\oplus r}}^d\simeq\bigoplus_{d_1+\cdots+d_r=d}\Av_{\mathbf1}^{d_1}\circ\cdots\circ\Av_{\mathbf1}^{d_r}.\]
所以，只需证明：对尖点 $\cF$ 及 $d>m(2g-2)$，
\[\Av_{\mathbf1}^d(\cF)=0.\]
最后一步是直接计算。在函数层面，它是 Jacquet–Godement 的陈述：$\GL_m$ 的尖点自守表示的标准 $L$ 函数，是次数至多 $m(2g-2)$ 的多项式；Braverman 的计算将它几何化。
\section{为何平均函子的像是尖点的}\label{sec:9.7}
令 $m=m_1+m_2$，并令
\[\CT:D(\Bun_m)\longrightarrow D(\Bun_{m_1}\times\Bun_{m_2})\]
为 Levi 子群为 $\GL_{m_1}\times\GL_{m_2}$ 的标准抛物子群的常数项函子。$\CT\circ\Av_E^d$ 上有一个过滤，其相伴分次部分为
\[\bigoplus_{d_1+d_2=d}(\Av_E^{d_1}\boxtimes\Av_E^{d_2})\circ\CT.\]
这是修改叠的纯几何过滤。由于对两个 Levi 因子都有 $\rank(E)>m_i$，当 $d>rm(2g-2)$ 时，较小秩的消失定理通过归纳消去每个分次部分：对每个 $d=d_1+d_2$，至少有一个 $d_i$ 超过 $rm_i(2g-2)$。故
\[\CT\circ\Av_E^d=0,\]
所以 $\Av_E^d$ 的像是尖点的。

这也迅速证明了上文用到的一个事实。若 $\cF\in D(\Bun_m)$ 是 $E$-Hecke 特征层，则如谱作用所预言，
\[\Av_E^d(\cF)\simeq\cF\otimes\RGamma(X^{(d)},(E^*\otimes E)^{(d)}).\]
应用 $\CT$ 得到
\[\CT(\Av_E^d(\cF))\simeq\CT(\cF)\otimes\RGamma(X^{(d)},(E^*\otimes E)^{(d)}).\]
另一方面，上述过滤将同一个对象表达为作用在 $\CT(\cF)$ 上的较小秩平均函子的和，因此对 $d\gg0$ 它消失。由于 $E^*\otimes E$ 包含平凡局部系统，所示上同调对适当的大 $d$ 非零。因此 $\CT(\cF)=0$，即 $\cF$ 是尖点的。
\newpage
\section{\texorpdfstring{$t$}{t} 正合性背后的商范畴}\label{sec:9.8}\index{Verdier 商|hyperpage}\index{t正合性@$t$ 正合性|hyperpage}\index{保守函子|hyperpage}
最初证明中的困难是 $t$-正合性。看似诱人的陈述“$\Av_E^1$ 是 $t$-正合的”是错误的。正确做法是构造一个 Verdier 商
\[D(\Bun_m)\twoheadrightarrow\widetilde D(\Bun_m),\]
具有三项性质：
\origdiagram{p052-square.pdf}{0.3370}
其中 $\widetilde{\Av}_E^1$ 是 $t$-正合的，该商继承反常 $t$-结构，而且复合
\[D(\Bun_m)_{\mathrm{cusp}}\longrightarrow D(\Bun_m)\longrightarrow\widetilde D(\Bun_m)\]
\srcpage{53}
是保守的。这是对 $\Av_E^1$ 那个错误正合性陈述的 Springer 型替代。Gaitsgory 最初构造 $\widetilde D$ 的方法在技术上相当复杂。

现代的解释是应将这个商取为温和部分：
\[\widetilde D(\Bun_G)=D(\Bun_G)^{\mathrm{temp}}.\]
几何 Ramanujan 定理说明，尖点对象是温和的\cite{Ber21}；Eisenstein／常数项的形式体系补上其余的非尖点部分。这为旧的商范畴提供了概念上更清楚的替代。

\begin{transnote}{背景直觉}
几何 Ramanujan 定理是经典 Ramanujan 猜想的几何版本：尖点自守表示的「特征值」都落在温度带内，等价地说，尖点对象在谱侧是温和的。经典对应是 $\GL_2$ 上尖点自守表示的 $L$ 函数在临界带上无极点（Ramanujan 界）；这里把它翻译成「尖点层属于 $\Shv_{\mathrm{Nilp}}$ 的温和子范畴」。它补上的正是旧的商范畴（尖点部分）缺掉的那一半：非尖点部分交给 Eisenstein／常数项形式体系。[依据：原书 p.53 本段；并见 [Ber21]。]
\end{transnote}
\endgroup
